\documentclass[aspectratio=169,11pt]{beamer}
\usepackage{fontspec}
\usepackage{listings}
\usepackage{xcolor}
\usepackage{booktabs}
\usepackage{array}

% ================= Palette (identique au livre Tome 1) =================
\definecolor{PyBleu}{HTML}{1A4A7A}
\definecolor{PyBleuMoyen}{HTML}{2D7AB3}
\definecolor{PyBleuClair}{HTML}{EAF1F8}
\definecolor{PyJaune}{HTML}{FFD43B}
\definecolor{PyNuit}{HTML}{142840}
\definecolor{PyTexte}{HTML}{1F2A37}
\definecolor{PyGris}{HTML}{5A6B7E}
\definecolor{PyRouge}{HTML}{B3442D}
\definecolor{PyVert}{HTML}{2E7D5B}
\definecolor{CodeTexte}{HTML}{E8EEF6}

% ================= Code Python =================
\lstdefinestyle{py}{
  language=Python,
  keywordstyle=\color{PyJaune}\bfseries,
  stringstyle=\color{PyJaune},
  commentstyle=\color{PyGris}\itshape,
  showstringspaces=false,
  breaklines=true,
  backgroundcolor=\color{PyNuit},
  basicstyle=\ttfamily\footnotesize\color{CodeTexte},
  columns=fullflexible,
  keepspaces=true,
  xleftmargin=6pt,
  framexleftmargin=6pt,
}
\lstdefinestyle{pysmall}{
  language=Python,
  keywordstyle=\color{PyJaune}\bfseries,
  stringstyle=\color{PyJaune},
  commentstyle=\color{PyGris}\itshape,
  showstringspaces=false,
  breaklines=true,
  backgroundcolor=\color{PyNuit},
  basicstyle=\ttfamily\scriptsize\color{CodeTexte},
  columns=fullflexible,
  keepspaces=true,
  xleftmargin=5pt,
  framexleftmargin=5pt,
}

% ================= Titres de diapositives =================
\setbeamercolor{structure}{fg=PyBleu}
\setbeamerfont{frametitle}{series=\bfseries, size=\Large}
\setbeamercolor{frametitle}{bg=PyBleu, fg=white}
\setbeamertemplate{frametitle}{%
  \nointerlineskip
  \begin{beamercolorbox}[wd=\textwidth, ht=2.1ex, dp=0.9ex,
                          leftskip=0.35cm, rightskip=0.35cm]{frametitle}%
    \usebeamerfont{frametitle}\insertframetitle%
  \end{beamercolorbox}%
  \vspace*{-0.4pt}%
  {\color{PyJaune}\rule{\textwidth}{1.6pt}}%
}

% ================= Pied de page =================
\setbeamerfont{footline}{size=\tiny}
\setbeamercolor{footline}{bg=white, fg=PyGris}
\newcommand{\PiedDePage}{Python en 24 heures \textbar\ Séance 03 : Variables et types}
\setbeamertemplate{footline}{%
  \leavevmode%
  \hbox{%
    \begin{beamercolorbox}[wd=\paperwidth, ht=2.2ex, dp=1.1ex,
                            leftskip=0.35cm, rightskip=0.35cm]{footline}%
      \usebeamerfont{footline}\PiedDePage
      \hfill\insertframenumber\,/\,\inserttotalframenumber%
    \end{beamercolorbox}}%
}

\setbeamertemplate{navigation symbols}{}
\setbeamertemplate{itemize item}{\color{PyBleuMoyen}$\blacktriangleright$}
\setbeamertemplate{itemize subitem}{\color{PyBleuMoyen}\textbullet}

% ================= Blocs pedagogiques =================
\setbeamercolor{block title}{bg=PyBleu, fg=white}
\setbeamercolor{block body}{bg=PyBleuClair, fg=PyTexte}
\setbeamercolor{block title alerted}{bg=PyRouge, fg=white}
\setbeamercolor{block body alerted}{bg=PyRouge!9, fg=PyTexte}
\setbeamercolor{block title example}{bg=PyJaune, fg=PyNuit}
\setbeamercolor{block body example}{bg=PyJaune!14, fg=PyTexte}

% ================= Metadonnees =================
\hypersetup{
  pdftitle={Python en 24 heures : Séance 03 : Variables et types},
  pdfauthor={HOUNKPE Mawulikplimi Roland et DJABON Ounimborbitibou},
  pdfsubject={De l'algorithmique à la programmation orientée objet},
  pdfkeywords={Python, algorithmique, POO, formation, exercices corrigés}
}

\begin{document}

{%
\setbeamercolor{background canvas}{bg=PyNuit}
\begin{frame}[plain]
\vspace{10mm}
{\color{PyJaune}\small\bfseries FORMATION « PYTHON EN 24 HEURES » \quad\textbar\quad SÉANCE 03\par}
\vspace{4mm}
{\color{white}\Huge\bfseries Premiers pas en Python : variables et types\par}
\vspace{4mm}
{\color{PyBleuClair}\large Chapitres 5 et 6 du Tome 1\par}
\vspace{7mm}
{\color{PyJaune}\rule{34mm}{2.2pt}\par}
\vspace{7mm}
{\color{white}\normalsize 1 h de cours (présentation Beamer) \quad$\bullet$\quad 1 h de pratique (notebook Jupyter)\par}
\vspace{2.5mm}
{\color{PyBleuClair}\normalsize HOUNKPE Mawulikplimi Roland et DJABON Ounimborbitibou\par}
\end{frame}
}
\begin{frame}{Objectifs de la séance}
\begin{itemize}
  \item Exécuter du code Python de trois façons : interpréteur, script, notebook
  \item Créer des variables et respecter les règles de nommage
  \item Manipuler les quatre types de base : int, float, str et bool
  \item Convertir proprement avec int(), float() et str()
  \item Formater des messages lisibles avec les f-strings
\end{itemize}
\end{frame}
\begin{frame}{Rappel de la séance 02}
Nous quittons le pseudo-code : les algorithmes des deux premières séances vont maintenant s'écrire en vrai Python.

\begin{itemize}
  \item Boucles POUR et TANT QUE : répéter sans recopier le code
  \item Compteur et accumulateur : les deux schémas fondamentaux
  \item Champion courant : minimum et maximum d'une série
  \item Tri par sélection : boucles imbriquées et échange
  \item Pièges : boucle infinie et décalage de un
\end{itemize}
\end{frame}
\begin{frame}{Trois façons d'exécuter du Python}
\begin{columns}[t]
\begin{column}{0.315\textwidth}
{\bfseries\color{PyBleu}Interpréteur}

\begin{itemize}
  \setlength{\itemsep}{1.5pt}\setlength{\parskip}{0pt}
  \item Commande `python -V` pour vérifier l'installation
  \item Tapez `python` : chaque ligne s'exécute aussitôt
  \item Idéal pour essayer une expression
\end{itemize}
\end{column}
\hspace{0.015\textwidth}
\begin{column}{0.315\textwidth}
{\bfseries\color{PyBleu}Script}

\begin{itemize}
  \setlength{\itemsep}{1.5pt}\setlength{\parskip}{0pt}
  \item Code dans un fichier `.py`
  \item Lancement : `python monfichier.py`
  \item Idéal pour un programme complet
\end{itemize}
\end{column}
\hspace{0.015\textwidth}
\begin{column}{0.315\textwidth}
{\bfseries\color{PyBleu}Notebook}

\begin{itemize}
  \setlength{\itemsep}{1.5pt}\setlength{\parskip}{0pt}
  \item Fichier `.ipynb` en cellules
  \item Chaque cellule se lance seule, le résultat reste affiché
  \item Notre format de TP : mélange d'explications et de code
\end{itemize}
\end{column}
\end{columns}
\end{frame}

\begin{frame}[fragile]{Afficher et commenter}

\begin{lstlisting}[style=py]
# Un programme Python se lit de haut en bas
print("Bonjour !")
print(3 + 4)
print("3 + 4 =", 3 + 4)
\end{lstlisting}
\vspace{0.5pt}

{\small\color{PyGris}\itshape Le dièse ouvre un commentaire : Python l'ignore totalement, il est là pour le lecteur humain.}

\end{frame}

\begin{frame}[fragile]{Variables et nommage}
\begin{columns}[t]
\begin{column}{0.42\textwidth}
\begin{itemize}
  \setlength{\itemsep}{1.5pt}\setlength{\parskip}{0pt}
  \item Affectation : `age = 25` crée la boîte et la remplit
  \item On peut relire et \textbf{réaffecter} : `age = age + 1`
  \item Noms en minuscules avec des tirets bas : `prix\_total`
  \item Interdits : espaces, accents, mots réservés (`if`, `for`...)
  \item Un nom clair vaut dix commentaires
\end{itemize}
\end{column}
\hspace{0.02\textwidth}
\begin{column}{0.55\textwidth}
\begin{lstlisting}[style=pysmall]
prix_total = 1500
remise = prix_total * 0.1
prix_total = prix_total - remise
print(prix_total)
\end{lstlisting}
\end{column}
\end{columns}
\end{frame}

\begin{frame}{Les quatre types de base}
\renewcommand{\arraystretch}{1.3}\small
\begin{center}
\begin{tabular}{lll}
\toprule
\textbf{\color{PyBleu}Type} & \textbf{\color{PyBleu}Contient} & \textbf{\color{PyBleu}Exemples} \\
\midrule
int & entiers, sans limite de taille & 0, 42, -7 \\
float & nombres à virgule flottante & 3.14, -0.5, 1e6 \\
str & texte entre guillemets & « bonjour », 'A', "" \\
bool & valeurs logiques & True, False \\
\bottomrule
\end{tabular}
\end{center}

\vspace{4pt}

{\small\color{PyGris}La fonction type(x) révèle le type d'une expression : essayez type(3), type(3.0), type(« 3 »).}
\end{frame}

\begin{frame}[fragile]{Les opérateurs arithmétiques}

\begin{lstlisting}[style=py]
print(7 + 2, 7 - 2, 7 * 2)   # 9 5 14
print(7 / 2)     # 3.5  division classique, toujours float
print(7 // 2)    # 3    division entière (partie entière)
print(7 % 2)     # 1    reste de la division entière
print(7 ** 2)    # 49   puissance
\end{lstlisting}
\vspace{0.5pt}

{\small\color{PyGris}\itshape Le quotient et le reste forment un duo : 7 = 2 fois 3 plus 1. `\%` sert pour la parité, les rotations, le découpage du temps.}

\end{frame}

\begin{frame}[fragile]{Priorités et parenthèses}
\begin{columns}[t]
\begin{column}{0.42\textwidth}
\begin{itemize}
  \setlength{\itemsep}{1.5pt}\setlength{\parskip}{0pt}
  \item Puissance avant multiplication et division
  \item Multiplication et division avant addition et soustraction
  \item En cas de doute : mettez des parenthèses
  \item Les parenthèses ne coûtent rien et évitent les erreurs
\end{itemize}
\end{column}
\hspace{0.02\textwidth}
\begin{column}{0.55\textwidth}
\begin{lstlisting}[style=pysmall]
moyenne = (12 + 15 + 9) / 3
print(moyenne)          # 12.0

sans_parentheses = 12 + 15 + 9 / 3
print(sans_parentheses)  # 30.0 : faux !
\end{lstlisting}
\end{column}
\end{columns}
\end{frame}

\begin{frame}{Conversions de types}
\begin{itemize}
  \item `int("12")` transforme un texte en entier ; `float("3.5")` en nombre réel
  \item `str(42)` transforme un nombre en texte : obligatoire pour coller des chaînes autrement
  \item `int(3.9)` \textbf{tronque} : résultat 3, sans arrondi
  \item Une conversion impossible lève `ValueError` : `int("douze")`
  \item Règle d'or : convertissez à la frontière, entre le monde du texte et celui des nombres
\end{itemize}
\end{frame}

\begin{frame}[fragile]{Les f-strings : écrire proprement}

\begin{lstlisting}[style=py]
prenom = "Awa"
note = 15.5
print(f"{prenom} a obtenu {note} sur 20")
print(f"Moyenne arrondie : {note:.1f}")
print(f"Sur 3 caractères : |{note:>8.2f}|")
\end{lstlisting}
\vspace{0.5pt}

{\small\color{PyGris}\itshape Le f devant les guillemets active l'interpolation : les expressions entre accolades sont évaluées puis insérées.}

\end{frame}

\begin{frame}[fragile]{Les booléens en action}

\begin{lstlisting}[style=py]
majeur = 19 >= 18
print(majeur)            # True
print(type(majeur))      # <class 'bool'>
print(2 == 2, 2 != 2)    # True False
print(True + True)       # 2 : True vaut 1 en calcul
\end{lstlisting}
\vspace{0.5pt}

{\small\color{PyGris}\itshape Toute comparaison produit un bool. Les connecteurs and, or et not s'appliquent aux bool, comme ET, OU, NON du pseudo-code.}

\end{frame}

\begin{frame}{Attention}
\begin{alertblock}{int n'est pas float n'est pas str}
`5` et `5.0` et `"5"` sont trois valeurs différentes : le premier est un int, le deuxième un float, la troisième une str. Additionner une str et un int lève `TypeError`. Le piège classique : lire une saisie avec `input()`, qui renvoie toujours une str, et oublier de la convertir avant de calculer.
\end{alertblock}
\end{frame}

\begin{frame}{Astuce}
\begin{exampleblock}{Explorer en direct}
Gardez un interpréteur ouvert pendant tout le TP. Une expression vous semble étrange ? Tapez la et observez. `type()`, `int()`, `round()` et les f-strings s'apprennent en cinq essais, pas en cinq pages de lecture.
\end{exampleblock}
\end{frame}

\begin{frame}{À retenir}
\begin{block}{L'essentiel de la séance 03}
Une variable porte un nom clair et garde une valeur typée : int, float, str ou bool. Les opérateurs `//` et `\%` découpent les entiers, `/` produit toujours un float. Les conversions se font explicitement avec int(), float() et str(). Les f-strings assemblent textes et nombres sans efforts ni erreurs.
\end{block}
\end{frame}

\begin{frame}{Travaux pratiques (1 h) : Variables, types et formatage}
Cinq exercices pour installer vos réflexes Python : affectations, calculs arithmétiques, conversions et f-strings. Exécutez chaque cellule après l'avoir complétée.

\begin{itemize}
  \setlength{\itemsep}{1.5pt}\setlength{\parskip}{0pt}
  \item \textbf{Exercice 1 (Basique)} : Carte de visite en f-string
  \item \textbf{Exercice 2 (Basique)} : Convertisseur secondes vers h:m:s
  \item \textbf{Exercice 3 (Intermédiaire)} : Prix TTC avec arrondi
  \item \textbf{Exercice 4 (Intermédiaire)} : Échange de variables et alignement
  \item \textbf{Exercice 5 (Défi)} : Rendre la monnaie en FCFA
\end{itemize}
\vspace{4pt}

{\small\color{PyGris}Ouvrez le notebook « Seance\_03\_variables\_et\_types.ipynb » pour commencer.}
\end{frame}
\begin{frame}{Pour la prochaine séance}
\begin{exampleblock}{Conseil de travail}
Séance 04 : conditions et boucles en Python. Le pseudo-code des séances 01 et 02 va prendre sa forme définitive : if, elif, else, while, for, range, break et continue. Refaites l'exercice de la monnaie en expliquant chaque ligne à voix haute : c'est le meilleur entraînement.
\end{exampleblock}
\end{frame}

\end{document}
